\section{Regularity and a bound on the entire unit circle}
Let $D=z\,d/dz$. Since $d/d\theta=iD$ on a circle, the fourth
angular derivative is exactly $D^4$. We will prove
\begin{equation}\label{eq:C4}
 \left(\frac1{2\pi}\int_0^{2\pi}|D^4H(e^{i\theta})|^2\dd\theta\right)^{1/2}
 <C_*:=24744.587741.
\end{equation}
Two different estimates for the mixed derivatives of $F$ are used on
closed angular intervals. Before giving the numerical estimates, we
establish the boundary regularity needed to interpret this integral and
to pass from it to coefficient bounds.

\subsection{Gaussian smoothing in one coordinate}
For $0<r<1$, put $v=1-r$ and
\[
 G_r(w,x)=\E_Z f(\sqrt r\,w+\sqrt{1-r}\,Z,x)
 =\erf\left(\frac{\sqrt r\,w+L(x)}{\sqrt{2v}}\right).
\]
Here $Z$ is an independent standard real Gaussian. The generating
function or (\ref{eq:hermite-correlation}) shows that $G_r$ has Hermite
coefficient $r^{j/2}A_{jk}$. For nonnegative $p,q$ with $p+q\geq1$ define
\begin{equation}\label{eq:Spq}
 S_{pq}(r)=r^{-p}\norm{\partial_w^p\partial_x^qG_r}_{L^2(\gamma^2)}^2
 =\sum_{j,k}\fall{j}{p}\fall{k}{q}r^{j-p}A_{jk}^2,
\end{equation}
where $(j)_p=j!/(j-p)!$ if $j\geq p$ and is zero otherwise.
To justify this identity without assuming differentiability of $f$, note
that every fixed derivative of the explicit smooth function $G_r$ is a
polynomial in $w,x$ times a Gaussian in $\sqrt r\,w+L(x)$. It and all
lower derivatives have at most polynomial growth and belong to
$L^2(\gamma^2)$. Gaussian integration by parts therefore has vanishing
boundary terms. It identifies the Hermite coefficients of the derivative
as the shifted coefficients of $G_r$ multiplied by
$\sqrt{(j)_p(k)_q}$. Completeness and Parseval prove
(\ref{eq:Spq}), including finiteness.

Termwise differentiation of (\ref{eq:F}) now gives
\begin{equation}\label{eq:smoothing-bound}
 |\partial_\alpha^p\partial_\beta^qF(\alpha,\beta)|
 \leq\frac\pi2 S_{pq}(r),\qquad |\alpha|\leq r,\quad |\beta|\leq1.
\end{equation}
The differentiated series converges absolutely and uniformly on this
set, by (\ref{eq:Spq}); terms with insufficient indices vanish. Taking
$r=(1+\rho_0)/2$ and using $|P(z)|\leq\rho_0$ on $|z|\leq1$ proves
that every mixed derivative needed through order four is continuous
after composition on the closed disk. The chain rule consequently gives
a continuous boundary value for $D^4H$, with uniform convergence of
$D^4H(se^{i\theta})$ as $s\uparrow1$. This includes $z=1$ and $z=-1$.
This qualitative argument requires no numerical evaluation at the
parameter $(1+\rho_0)/2$.

We next give an explicit finite calculation of $S_{pq}(r)$ for
$1\leq p+q\leq4$. Put $T=\sqrt r\,w+L(x)$ and
$\psi(T)=\erf(T/\sqrt{2v})$. For $d\geq1$,
\[
 \psi^{(d)}(T)=\sqrt{2/\pi}\,v^{-1/2}e^{-T^2/(2v)}U_d(T),
 \qquad U_d(T)=(-1)^{d-1}v^{-(d-1)/2}\He_{d-1}(T/\sqrt v).
\]
In particular,
\[
 U_1=1,\quad U_2=-T/v,\quad U_3=T^2/v^2-1/v,
 \quad U_4=-T^3/v^3+3T/v^2.
\]
The chain rule gives
\begin{equation}\label{eq:R}
 \partial_w^p\partial_x^qG_r
 =r^{p/2}\sqrt{2/\pi}\,v^{-1/2}e^{-T^2/(2v)}R_{pq}(x,T),
 \quad R_{pq}=\sum_k B_{qk}U_{p+k}.
\end{equation}
Here $B_{00}=1$; for $q>0$ the nonzero Bell polynomials are
\begin{align*}
 B_{11}&=L',\\
 B_{21}&=L'',&B_{22}&=(L')^2,\\
 B_{31}&=L''',&B_{32}&=3L'L'',&B_{33}&=(L')^3,\\
 B_{41}&=L'''',&B_{42}&=4L'L'''+3(L'')^2,&
 B_{43}&=6(L')^2L'',&B_{44}&=(L')^4.
\end{align*}
These identities follow by differentiating the previous row and
collecting derivatives of $\psi$; no asymptotic chain rule is involved.
Only $p+k\geq1$ occurs when $p+q\geq1$.

For any polynomial $Q_0$ and fixed real $x$, completing the square in the
$w$ integral proves
\begin{equation}\label{eq:tilt}
 \E_w[e^{-T^2/v}Q_0(T)]
 =\sqrt{\frac v{1+r}}e^{-L(x)^2/(1+r)}\E[Q_0(Y)],
 \quad Y\sim N\left(\frac{vL(x)}{1+r},\frac{rv}{1+r}\right).
\end{equation}
Indeed the original density has variance $r$ and mean $L(x)$ in the
$T$ coordinate; adding $T^2/v$ to its exponent gives the stated mean,
variance, and normalizing factor. The moments of $Y$, with mean $\mu$
and variance $\sigma^2$, are exactly
\begin{equation}\label{eq:tilted-moments}
 \E Y^d=\sum_{\ell=0}^{\lfloor d/2\rfloor}
 \frac{d!}{2^\ell\ell!(d-2\ell)!}\sigma^{2\ell}\mu^{d-2\ell},
\end{equation}
as follows by expanding $\exp(\mu t+\sigma^2t^2/2)$.
With the real polynomial
$Q_{pq}(x)=\E[R_{pq}(x,Y)^2]$, equations
(\ref{eq:Spq})--(\ref{eq:tilted-moments}) give
\begin{equation}\label{eq:S-integral}
 S_{pq}(r)=\frac{2}{\pi\sqrt{1-r^2}}
       \int_\R Q_{pq}(x)e^{-L(x)^2/(1+r)}\dd\gamma(x).
\end{equation}
This polynomial is even, is nonnegative on the real line, and has degree
at most 118 in the cases used here. Its coefficients need not be
nonnegative. Evenness also follows directly from simultaneous sign
change of $w,x$ in the squared derivative in (\ref{eq:R}).

There is a second way to construct exactly the same polynomial, used in
the accompanying implementation. For polynomials in $T$ with polynomial
coefficients in $x$, let
\[
 \mathcal T R=\partial_T R-TR/v,\qquad
 \mathcal X R=\partial_xR+L'(x)\mathcal T R.
\]
If $p>0$, start from $1$, apply $\mathcal T$ a total of $p-1$ times,
then $\mathcal X$ a total of $q$ times. If $p=0$, start from $L'$ and
apply $\mathcal X$ a total of $q-1$ times. Differentiating
$e^{-T^2/(2v)}R$ verifies these rules directly. Squaring the result and
using (\ref{eq:tilted-moments}) gives $Q_{pq}$, providing a finite
construction independent of the Bell polynomial listing.

\subsection{The complete error for the smoothing integrals}
Every finite integral in (\ref{eq:S-integral}) is evaluated on $[0,16]$
as the sum of the 16 integrals on $[j,j+1]$, then doubled. The complex
integrand supplied to the interval integrator is the entire function
\begin{equation}\label{eq:entire-S}
 J(z)=\frac{2Q_{pq}(z)}{\pi\sqrt{1-r^2}\sqrt{2\pi}}
       \exp\left(-\frac{L(z)^2}{1+r}-\frac{z^2}{2}\right).
\end{equation}
In particular, the square in this formula is the analytic square, not a
complex absolute square. Enclosing polynomial coefficients enclose the
exact algebraic coefficients of $L$ and $Q_{pq}$.

The finite-interval calculation uses the certified complex integration
algorithm in Arb/FLINT~\cite{arb,integration,flintcalc}. Its adaptive
Gauss--Legendre rule encloses the integrand on a full complex ellipse;
it accepts a quadrature result only with a rigorous remainder. For an
ellipse of parameter $\rho>1$, maximum bound $M$, and $s$ nodes, the
bound used by that algorithm on $[-1,1]$ is
\begin{equation}\label{eq:library-error}
 \frac{64M}{15(\rho-1)\rho^{2s-1}}.
\end{equation}
It is the analytic Gauss--Legendre remainder bound implemented in the
cited algorithm; scaling by the interval half-width gives the bound on
a subinterval. A fully elementary, more conservative remainder
$8M\rho^{-2s}/(1-\rho^{-1})$ follows from the Chebyshev proof above and
also explains why such enclosing quadratures converge here. Subdivision
adds enclosing integrals and their errors, rather than extrapolating
point samples. Direct integration enclosures on small intervals are
also permitted by the algorithm. All calls use 384-bit arithmetic,
relative tolerance $10^{-18}$, absolute tolerance $10^{-20}$, and an
evaluation limit of $10^6$; these tolerances request accuracy and are
not themselves claimed as error bounds. Only the returned finite
enclosures are used. Entire analyticity of (\ref{eq:entire-S}) validates
every analyticity query and every ellipse bound required by the
integrator. The mathematical dependence here is the proved enclosing
quadrature algorithm and the inclusion property of ball arithmetic,
not an uncertified numerical integration routine.

The infinite remainder is bounded separately. Write
$L(x)=\sum_{k=0}^9\lambda_k x^k$. The coefficient enclosures prove
\begin{equation}\label{eq:Ltail}
 \sum_{k<9}|\lambda_k|16^{k-9}<\frac{|\lambda_9|}{2},
 \qquad c:=10^{-6}<\frac{|\lambda_9|}{2}.
\end{equation}
Hence $|L(x)|\geq c|x|^9$ for $|x|\geq16$. If
$m=\deg Q_{pq}$, the further checked inequality
\[
 \frac{18c^2 16^{18}}{1+r}>m
\]
shows that $x^ke^{-c^2x^{18}/(1+r)}$ decreases for $x\geq16$ and
every $0\leq k\leq m$. Therefore the entire two-sided omitted integral
in (\ref{eq:S-integral}) is at most
\begin{equation}\label{eq:S-tail}
 \frac{2}{\pi\sqrt{1-r^2}}
   \left(\sum_{k=0}^m|[x^k]Q_{pq}|16^k\right)
       e^{-c^2 16^{18}/(1+r)}<2^{-100}.
\end{equation}
In deriving this bound we used only that the Gaussian measure of
$\{|x|>16\}$ is at most one, so both tails are included. A symmetric
interval of this radius is added to twice the finite real integral.
All 266 required integral records contain the finite integral enclosure,
this tail bound, and the positive margins in (\ref{eq:Ltail}) and the
monotonicity test. These records are computed from an empty integral
cache.

\subsection{Mixed derivative bounds from complex Gaussian densities}
For $z\in\C$ with $|\Re z|<1$, the analytic continuation of a correlated
Gaussian density, after the real orthogonal change of variables
$U=(x+y)/\sqrt2$, $V=(x-y)/\sqrt2$, is
\[
 k_z(U,V)=\frac{1}{2\pi\sqrt{1+z}\sqrt{1-z}}
       \exp\left(-\frac{U^2}{2(1+z)}-\frac{V^2}{2(1-z)}\right).
\]
Each square root is its analytic branch in the right half-plane.
On compact subsets of the strip the real parts of the inverse variances
are positive and uniformly bounded below. Thus this density and every
fixed parameter derivative have an integrable Gaussian majorant times
a polynomial. Differentiation of its integral against bounded real
functions is justified by dominated convergence and defines a
holomorphic function.

Let $d_0(z)=\int_{\R^2}|k_z|$. Completing two real Gaussian integrals
gives
\begin{equation}\label{eq:density-mass}
 d_0(z)=\sqrt{\frac{|1+z|\,|1-z|}{1-(\Re z)^2}},\quad
 v_U=\frac{|1+z|^2}{1+\Re z},\quad
 v_V=\frac{|1-z|^2}{1-\Re z}.
\end{equation}
The probability density $|k_z|/d_0(z)$ is that of independent centered
real Gaussians with variances $v_U,v_V$. To compute derivatives write
$\partial_z^p k_z=k_z T_p(U,V;z)$. One explicit formula is
\begin{align}
 R_j(v,u)&=\sum_{\ell=0}^j
 \frac{(2j)!(-1)^{j-\ell}}{2^{2j-\ell}(j-\ell)!(2\ell)!}
                       v^{-j-\ell}u^{2\ell},\nonumber\\
 T_p&=\sum_{j=0}^p\binom pj(-1)^{p-j}
                          R_j(1+z,U)R_{p-j}(1-z,V).\label{eq:density-poly}
\end{align}
For a one-dimensional Gaussian density, the heat equation
$\partial_v=\tfrac12\partial_u^2$ and the Hermite formula for its
spatial derivatives prove $R_j$. The product rule and
$\partial_z(1-z)=-1$ then prove (\ref{eq:density-poly}).
Cauchy--Schwarz under the probability law just described gives
\begin{equation}\label{eq:density-bound}
 \norm{\partial_z^pk_z}_{L^1(\R^2)}
 \leq d_p(z):=d_0(z)\sqrt{\E|T_p(U,V;z)|^2}.
\end{equation}
All expectations here are finite polynomial Gaussian moments, without
any integral truncation.

For an alternative exact construction, the $k$th logarithmic derivative
$\ell_k=\partial_z^k\log k_z$ has constant, $U^2$, and $V^2$
coefficients respectively
\begin{align*}
 &\frac{(-1)^k(k-1)!}{2(1+z)^k}
           +\frac{(k-1)!}{2(1-z)^k},\\
 &\frac{(-1)^{k+1}k!}{2(1+z)^{k+1}},\qquad
           -\frac{k!}{2(1-z)^{k+1}}.
\end{align*}
Starting with $T_0=1$, the finite recurrence
\[
 T_m=\sum_{k=1}^m\binom{m-1}{k-1}\ell_kT_{m-k}
\]
is obtained by differentiating an exponential generating function for
$k_{z+t}/k_z$. It is the construction used in the independent
implementation. To evaluate its squared norm without a cancellation
assumption, rescale $U=\sqrt{v_U}\,u$, $V=\sqrt{v_V}\,v$, and use
\[
 u^{2k}=\sum_{i=0}^k
 \frac{(2k)!}{2^{k-i}(k-i)!(2i)!}\He_{2i}(u).
\]
This identity follows by comparing coefficients in the Hermite
generating function multiplied by $e^{t^2/2}$. If the resulting
polynomial is $\sum_{i,j}c_{ij}\He_{2i}(u)\He_{2j}(v)$, then
\[
 \E|T_p|^2=\sum_{i,j}|c_{ij}|^2(2i)!(2j)!.
\]
Thus all complex conjugations occur only in a finite nonnegative square
norm, not in the analytic density or its derivative.

Using independent copies of the correlated density for the $W$ and $X$
coordinates, boundedness $|fg|\leq1$ gives
\begin{equation}\label{eq:mixed-density}
 |\partial_\alpha^p\partial_\beta^qF(\alpha,\beta)|
 \leq\frac\pi2d_p(\alpha)d_q(\beta),
 \qquad |\Re\alpha|,|\Re\beta|<1.
\end{equation}
To identify this analytic function with (\ref{eq:F}), first note their
equality on the real open square of correlations. Fixing a real second
coordinate and using the one-variable identity theorem, then fixing a
complex first coordinate and repeating, proves equality in the open
bidisk. At boundary points within the strip, the derivatives agree with
the continuous limits established by (\ref{eq:Spq}). This proves that
(\ref{eq:mixed-density}) bounds the same derivatives used by the
coefficient argument.

\subsection{Angular chain rule and complete interval coverage}
For $s\geq0$ write
\[
 D^sH(z)=\sum_{p,q}c^{(s)}_{pq}(z)
               (\partial_\alpha^p\partial_\beta^qF)(P(z),z).
\]
These are exact rational polynomial coefficients. For example they are
given by the finite formula
\begin{equation}\label{eq:angular-jets}
 c^{(s)}_{pq}(z)=\frac{s!}{p!q!}[u^s]
             (P(ze^u)-P(z))^p(ze^u-z)^q.
\end{equation}
Taylor's formula for $F(P(ze^u),ze^u)$ proves this identity in the open
disk. It can also be computed recursively, with out-of-range indices
zero, from
\begin{equation}\label{eq:angular-rec}
 c^{(0)}_{00}=1,\qquad
 c^{(s+1)}_{pq}=D c^{(s)}_{pq}+(DP)c^{(s)}_{p-1,q}+z c^{(s)}_{p,q-1}.
\end{equation}
For $s=4$ the nonzero indices are precisely
$p,q\geq0$, $1\leq p+q\leq4$, a set of 14 elements. All these
polynomials are computed over $\mathbb Q$, so the finite chain rule
contains no numerical differentiation error.

Divide the first quadrant into the 512 closed intervals
\[
 I_j=[j\pi/1024,(j+1)\pi/1024],\qquad 0\leq j<512.
\]
Interval sine and cosine of $kI_j$ enclose $z^k=e^{ik\theta}$ for
every $\theta\in I_j$. This yields enclosures for $P(z)$ and all
$c^{(4)}_{pq}(z)$ on each whole interval. A smoothing bound is available
whenever the enclosure proves $|P(z)|<r$ for one of
$r=78/100,79/100,\ldots,97/100$. A density bound is available when
both $|\Re P(z)|<1$ and $|\Re z|<1$ are proved by the same interval
enclosures. Every panel has at least one available estimate. Where
both are available, their smaller upper bound may be used separately
for each $(p,q)$. The intervals touching $z=1$ are covered by smoothing;
the density estimate is not asserted at its singular endpoint.

Let $U_{j,pq}$ be such an upper bound for
$(2/\pi)|\partial_\alpha^p\partial_\beta^qF|$, and let $C_{j,pq}$
bound $|c^{(4)}_{pq}|$. Then
\begin{equation}\label{eq:panel-majorant}
 |D^4H(e^{i\theta})|\leq M_j:=\frac\pi2
             \sum_{1\leq p+q\leq4}C_{j,pq}U_{j,pq}
             \quad(\theta\in I_j).
\end{equation}
The circle certificate records all $512\cdot14=7168$ terms, each
chosen method, and its domain inequalities. Exactly 266 smoothing
integrals are required: all 14 indices at each of the 19 parameters
$79/100,\ldots,97/100$. These are the integrals bounded in
(\ref{eq:S-integral})--(\ref{eq:S-tail}).

Real coefficients and oddness imply
$D^4H(\bar z)=\overline{D^4H(z)}$ and $D^4H(-z)=-D^4H(z)$.
The four quadrants therefore give equal contributions to its squared
norm. Round each $M_j$ upwards to $\widehat M_j=k_j/10^8$ with
integer $k_j$. The complete list of these 512 rational upper bounds is
included in the certificate. Summing their exact rational squares gives
\begin{equation}\label{eq:rational-square}
 \frac1{2\pi}\int_0^{2\pi}|D^4H(e^{i\theta})|^2\dd\theta
 \leq\frac1{512}\sum_{j=0}^{511}\widehat M_j^2
 =Q:=\frac{1567474233421402382277578179}{2560000000000000000}.
\end{equation}
The explicit strict rational comparison
\begin{equation}\label{eq:Cmargin}
 C_*^2-Q=\frac{107039174409781821}{2560000000000000000}>0
\end{equation}
proves (\ref{eq:C4}). The bounds concern complete angular intervals,
not samples of the circle.

\section{The coefficient tail and the exact numerical conclusion}
For $0<s<1$, ordinary Fourier Parseval applied to the power series gives
\[
 \frac1{2\pi}\int_0^{2\pi}|D^4H(se^{i\theta})|^2\dd\theta
       =\sum_{n\geq0}n^8|b_n|^2s^{2n}.
\]
The continuous boundary extension proved above permits passage to the
limit on the left by uniform convergence. Monotone convergence of the
nonnegative series gives the limit on the right. Hence
\begin{equation}\label{eq:parseval-boundary}
 \sum_{n\geq0}n^8|b_n|^2
 =\frac1{2\pi}\int_0^{2\pi}|D^4H(e^{i\theta})|^2\dd\theta<C_*^2.
\end{equation}
For every positive odd $N$, Cauchy--Schwarz and the decreasing function
$(N+2x)^{-8}$ show that
\begin{align}
 \sum_{\substack{n>N\\n\ {\rm odd}}}|b_n|
 &\leq\left(\sum_{n>N}n^8|b_n|^2\right)^{1/2}
        \left(\sum_{k=1}^{\infty}(N+2k)^{-8}\right)^{1/2}\nonumber\\
 &<\frac{C_*}{\sqrt{14N^7}},\label{eq:tail-final}
\end{align}
because the latter sum is at most
$\int_0^\infty(N+2x)^{-8}\dd x=1/(14N^7)$.
For $N=501$ the first omitted coefficient is $b_{503}$.
